% !TEX root = ../thesis.tex
 
\chapter{Záver}\label{ch:summary}
 
Práca sa zamerala na návrh a realizáciu AI asistenta pre oblasť energetiky v Európe. Využité boli malé jazykové modely a výpočtovo úsporné metódy doučenia. Hlavným cieľom bolo overiť, či je možné dosiahnuť prakticky použiteľné doménové odpovede aj bez veľkých výpočtových klastrov.
 
\section*{Zhrnutie teoretickej časti}
 
V analytickej časti bol spracovaný vývoj jazykových modelov. Osobitná pozornosť bola venovaná parameter-efektívnym fine-tuning technikám. Boli analyzované tri malé modely: Phi-3 Mini (3.8B), Gemma 3 4B a Qwen3 4B. Pre každý z nich boli posúdené architektonické vlastnosti, licenčné podmienky a kompatibilita s dostupnými knižnicami. Na základe tejto analýzy bol zvolený postup spájajúci Instruct inicializáciu, SFT, QLoRA a doplnkovú DPO kalibráciu. Táto kombinácia je vhodná pre prostredia s obmedzeným hardvérovým vybavením.
 
\section*{Zhrnutie praktickej časti}
 
V praktickej časti bol implementovaný celý pipeline doménovej adaptácie. Tréningový dataset bol zostavený z odborných zdrojov pokrývajúcich európsku energetiku. Dataset obsahuje 1159 vzoriek pre SFT a 1126 vzoriek pre DPO. Každý model bol trénovaný s vlastnou konfiguráciou hyperparametrov. Zdieľaná dátová pipeline zaistila jednotné predspracovanie naprieč všetkými modelmi.
 
Praktickým výstupom práce je reprodukovateľná tréningová pipeline. Jej súčasťou je aj funkčný konverzačný systém postavený na FastAPI. Systém umožňuje prepínanie modelov a generovanie štruktúrovaných odpovedí. Môže bežať na obmedzenom hardvéri bez špeciálneho výpočtového klastra.
 
\section*{Splnenie stanovených cieľov}
 
Stanovené ciele práce boli splnené nasledovne:
 
\begin{itemize}
    \item \textbf{Príprava datasetu:} bol zostavený doménový dataset z energetických zdrojov vhodný pre inštrukčný tréning, vrátane preferenčných párov pre DPO.
    \item \textbf{Výber a porovnanie modelov:} boli identifikované a porovnané tri SLM modely podľa veľkosti, architektúry a dostupnosti.
    \item \textbf{Doučenie modelov:} všetky tri modely boli doučené pomocou QLoRA na doménovom datasete. Výpočtové nároky boli zdokumentované.
    \item \textbf{Konverzačný asistent:} bol navrhnutý a implementovaný jednotný inferenčný systém vo forme REST API.
    \item \textbf{Vyhodnotenie:} schopnosti modelov boli zmerané pomocou G-Eval, perplexity, TTFT, TPS a VRAM profilovania.
    \item \textbf{Porovnanie s všeobecnými modelmi:} výsledky potvrdili, že malé doučené modely môžu predstavovať realistickú alternatívu k veľkým všeobecným modelom v doménovo špecifických scenároch.
\end{itemize}
 
\section*{Dosiahnuté výsledky}
 
Napriek obmedzeniu na 20~GB VRAM a datasetu s 1159 príkladmi boli dosiahnuté merateľné zlepšenia. Priemerné G-Eval skóre vzrástlo oproti base variantom o 3.77~\% (Phi-3 SFT) a 5.04~\% (Phi-3 DPO). Pri Gemma-3 a Qwen3 boli prínosy jemnejšie (0.40--1.25~\%), no konzistentné. Prejavili sa najmä v kategórii factual accuracy a v skracovaní odpovedí.
 
Z pohľadu kompromisu medzi kvalitou a pamäťovými nárokmi sa ako najvyváženejšia voľba ukázal model Qwen3 SFT. Dosahuje G-Eval skóre 9.53 v kategórii factual accuracy a patrí k pamäťovo najúspornejším 4B modelom. Ak je prioritou maximálna faktická správnosť, najlepší výsledok prináša Gemma-3 DPO (9.63). Pre scenáre s extrémne obmedzenou VRAM a nízkou latenciou je vhodný Phi-3 Base (TTFT 0.083~s, VRAM 858~MB).
 
Doménové doučenie tiež výrazne skrátilo priemerné odpovede. Všetky SFT/DPO varianty generujú odpovede o 50--70~\% kratšie než base modely. To je v súlade s cieľom stručného a štruktúrovaného výstupu vhodného pre rozhodovacie procesy.
 
\section*{Obmedzenia práce}
 
Práca má niekoľko obmedzení. Dataset 1159 príkladov je kompaktný. Pre robustnejšie závery by bol vhodný väčší a tematicky diverzifikovanejší dataset. Evaluácia bola realizovaná na 30 testovacích vzorkách bez štandardnej odchýlky. Výsledky preto treba interpretovať ako orientačné. G-Eval hodnotiteľ bol lokálny model Mistral. Iný hodnotiteľ mohol by priniesť odlišné skóre.
 
\section*{Možnosti ďalšieho výskumu}
 
Práca otvára niekoľko smerov pre pokračovanie. Prvým je rozšírenie datasetu o nové zdroje a jazyky. Väčší a diverzifikovanejší dataset by zlepšil schopnosť generalizácie. Druhým smerom je pridanie retrieval mechanizmu (RAG). Ten by umožnil modelu odkazovať na konkrétne dokumenty a citovať zdroje. Faktická spoľahlivosť by sa tým zvýšila. Tretím smerom je meranie robustnosti voči nejednoznačným otázkam a halucináciám. Štvrtým je porovnanie SFT a DPO variantov na širšej sade metrík vrátane ľudského hodnotenia. Piaty smer predstavuje nasadenie systému v reálnom prostredí a zber spätnej väzby od doménových expertov.
 
Tieto kroky by mohli ďalej zvýšiť praktickú využiteľnosť systému. Rovnako by mohli potvrdiť výsledky tejto práce na väčšom rozsahu experimentov.